\begin{table}[H]
\centering
\caption{Effect of a 27\% shock to extraction and refining with cross-country propagation, 2021--2023}
\label{tab:shock-oil-gas-chain-ea20}
\scriptsize
\setlength{\tabcolsep}{1.2pt}
\begin{tabular}{lrrrrrr}
\toprule
Country & \shortstack{Effect on average\\inflation -- Total} & \shortstack{Direct effect on\\average inflation} & \shortstack{Indirect effect on\\average inflation} & \shortstack{Effect on\\inflation gap} & \shortstack{Direct effect on\\inflation gap} & \shortstack{Indirect effect on\\inflation gap}
 \\
\midrule
Austria & 2.2 & 1.3 & 1.0 & 0.5 & 0.3 & 0.1
 \\
Belgium & 3.1 & 1.8 & 1.3 & 1.3 & 1.1 & 0.2
 \\
Cyprus & 2.7 & 2.2 & 0.5 & 1.4 & 1.2 & 0.2
 \\
Germany & 2.8 & 2.0 & 0.8 & 1.0 & 1.0 & 0.1
 \\
Estonia & 3.0 & 1.4 & 1.5 & 1.5 & 0.6 & 0.9
 \\
Greece & 2.4 & 1.9 & 0.6 & 1.4 & 1.3 & 0.1
 \\
Spain & 2.1 & 1.5 & 0.6 & 0.1 & 0.1 & 0.1
 \\
Finland & 1.9 & 1.4 & 0.6 & -0.7 & -0.7 & -0.0
 \\
France & 2.1 & 1.6 & 0.6 & 0.4 & 0.4 & -0.0
 \\
Croatia & 4.0 & 2.9 & 1.1 & 3.4 & 2.3 & 1.1
 \\
Ireland & 1.5 & 1.2 & 0.3 & 1.2 & 1.1 & 0.1
 \\
Italy & 3.5 & 2.3 & 1.1 & 1.2 & 1.0 & 0.2
 \\
Lithuania & 2.8 & 2.3 & 0.5 & 0.6 & 0.6 & -0.0
 \\
Luxembourg & 3.4 & 3.1 & 0.3 & 0.3 & 0.2 & 0.0
 \\
Latvia & 2.9 & 2.1 & 0.8 & 2.2 & 1.6 & 0.7
 \\
Malta & 1.5 & 1.2 & 0.3 & 0.1 & 0.1 & -0.0
 \\
Netherlands & 2.5 & 2.0 & 0.5 & 0.9 & 0.9 & -0.0
 \\
Portugal & 2.0 & 1.3 & 0.7 & 0.8 & 0.6 & 0.2
 \\
Slovenia & 3.0 & 2.6 & 0.4 & 1.2 & 1.2 & 0.1
 \\
Slovakia & 1.6 & 1.2 & 0.5 & 0.4 & 0.4 & 0.0
 \\
\midrule
\textbf{Euro area} & \textbf{2.6} & \textbf{1.8} & \textbf{0.8} & \textbf{0.8} & \textbf{0.7} & \textbf{0.1}
 \\
\bottomrule
\end{tabular}
\par\vspace{0.35em}\footnotesize\noindent Sources: Eurostat HICP, HBS, FIGARO, authors' calculations.
\par\vspace{0.25em}\footnotesize\noindent Notes. Entries are percentage-point effects over 2021--2023 of a uniform 27 percent shock to extraction and refining, calibrated as the simple average of the annual percentage changes in Brent crude oil and European TTF natural gas prices from 2021 to 2022 and from 2022 to 2023. The directly shocked FIGARO sectors are B (mining and quarrying, including NACE B06 crude petroleum and natural gas extraction) and C19 (coke and refined petroleum products). Sector D35 (electricity, gas, steam and air-conditioning supply) is not directly shocked. The shock vector is $s_{c,k}=0.27$ for $k\in\{B,C19\}$ and $s_{c,k}=0$ otherwise, for every country $c$. The price response is $\Delta p = (I - A^{\prime})^{-1}s$, where $A_{(o,i),(d,j)}=Z_{(o,i),(d,j)}/x_{d,j}$ retains bilateral intermediate flows from origin country $o$ and sector $i$ to destination country $d$ and sector $j$. The direct effect is $w_g^{\prime}Ms$ and the indirect effect is $w_g^{\prime}M[(I-A^{\prime})^{-1}-I]s$, where $M$ is the normalized NACE--COICOP bridge and $w_g$ contains household-group expenditure weights. The mapping assumes complete and immediate pass-through from FIGARO basic-price changes to the corresponding HICP categories; retail margins, taxes, subsidies, imported final goods, and incomplete pass-through are not separately modelled. The euro-area aggregate covers all 20 member countries. Imports from origins outside the EA20 production system enter production costs but are treated as exogenous and do not generate additional endogenous feedback. FIGARO flows are measured in current million euros at basic prices; the dataset used here does not provide separate import or export deflators. Income inflation gap = Q1 - Q5.
\end{table}
\vspace{0.6em}
